\section{Results}
\label{sec:results}

\subsection{Main Results}

Table~\ref{tab:main_v4} summarizes the primary benchmark. On the V4 main dataset with CultureMERT, the target-calibrated DCAS operating point outperforms the strongest hybrid baseline (BPR+LambdaMART) on all three manuscript-level metrics. Serendipity increases from 0.5558 to 0.8316, corresponding to a paired gain of +0.2757 with a 95\% bootstrap confidence interval (CI) of [0.2680, 0.2809]. At the same time, cultural calibration KL decreases from 2.0966 to 2.0296 ($\Delta=-0.0669$, 95\% CI [-0.0717, -0.0610]) and minority exposure at k rises from 0.2681 to 0.4023 ($\Delta=+0.1343$, 95\% CI [0.1275, 0.1404]). Relative to the strongest hybrid baseline, this corresponds to a 49.6\% gain in serendipity and a 50.1\% gain in minority exposure while also improving calibration.

This pattern matters because it rules out the simplest explanation that the observed gains come only from changing the embedding backbone or favoring one metric at the expense of the others. On the same V4 main protocol with Gemini, the target-calibrated operating point still improves over the strongest hybrid baseline, raising serendipity from 0.7884 to 0.8245 ($\Delta=+0.0361$, 95\% CI [0.0315, 0.0395]), lowering calibration KL from 2.3183 to 2.3104 ($\Delta=-0.0079$, 95\% CI [-0.0086, -0.0071]), and increasing minority exposure at k from 0.2746 to 0.3760 ($\Delta=+0.1014$, 95\% CI [0.0951, 0.1088]). The routeA_small sanity-check track shows the same direction even more clearly: with Gemini, serendipity increases from 0.7267 to 0.8642 ($\Delta=+0.1375$, 95\% CI [0.1282, 0.1499]), calibration KL drops from 1.5676 to 1.5502 ($\Delta=-0.0174$, 95\% CI [-0.0197, -0.0156]), and minority exposure at k rises from 0.1513 to 0.4997 ($\Delta=+0.3484$, 95\% CI [0.3301, 0.3677]).

Taken together, the main result is not that DCAS wins every operating point or every backbone indiscriminately. Rather, the consistent effect is that the downstream stack yields a better calibration-exposure-quality balance than baseline retrieval or hybrid reranking alone, which is the core empirical claim of the paper.

\subsection{Ablation Study}

We use two complementary ablations to test whether the framework's behavior can be reduced to a single component. The first is a classical module-removal study on the legacy \texttt{v2\_main} configuration with Gemini. Removing the pairwise constraints increases serendipity from 0.8225 to 0.8413 ($\Delta=+0.0188$, 95\% CI [0.0135, 0.0235]) but reduces minority exposure at k from 0.3629 to 0.3386 ($\Delta=-0.0243$, 95\% CI [-0.0308, -0.0172]), with no reliable change in calibration KL ($p=0.243$). Removing domain shaping produces a similar but smaller pattern: serendipity rises by +0.0087 (95\% CI [0.0037, 0.0135]) while minority exposure drops by -0.0158 (95\% CI [-0.0225, -0.0090]). By contrast, removing OT does not yield a stable global improvement in serendipity ($\Delta=+0.0014$, 95\% CI [-0.0002, 0.0031], $p=0.123$) and barely changes minority exposure ($p=0.056$). These results show that constraints and domain shaping are not decorative regularizers: they are the modules that keep the recommendation space from drifting toward a purely high-serendipity, low-exposure solution.

The second ablation treats calibration-aware reranking as an inference-time design decision rather than a fixed post-processing rule. On V4 main with CultureMERT, moving from the uncalibrated OT ranker to the target-calibrated operating point changes serendipity by $\Delta=-0.0263$ (95\% CI [-0.0285, -0.0247]), changes calibration KL by $\Delta=-0.0529$ (95\% CI [-0.0576, -0.0482]), and changes minority exposure at k by $\Delta=+0.1563$ (95\% CI [0.1529, 0.1597]). Moving further to the minority-oriented operating point continues to raise minority exposure from 0.4023 to 0.5303, at the cost of an additional 0.0034 drop in serendipity relative to the target-calibrated point. The same intervention is even cleaner on routeA_small with Gemini: the target-calibrated operating point does not significantly hurt serendipity relative to uncalibrated OT ($\Delta=+0.0036$, 95\% CI [-0.0025, 0.0100], $p=0.249$), but still changes calibration KL by $\Delta=-0.0221$ (95\% CI [-0.0234, -0.0211]) and changes minority exposure at k by $\Delta=+0.2574$ (95\% CI [0.2454, 0.2697]). This smooth frontier is why we describe calibration-aware reranking as a controllable operating-point mechanism rather than an ad hoc fairness tweak.

\subsection{Qualitative Analysis and Failure Modes}

% TODO: replace this subsection with item-level recommendation panels once case exports are available.
The most informative qualitative pattern is target-specific asymmetry rather than global failure. On V4 main with CultureMERT, the target-calibrated operating point improves serendipity for \texttt{modern\_english\_pop} ($\Delta=+0.0245$, 95\% CI [0.0196, 0.0298]) and leaves Turkey statistically ambiguous ($\Delta=+0.0031$, 95\% CI [-0.0012, 0.0071]), but reduces serendipity on China ($\Delta=-0.1007$, 95\% CI [-0.1090, -0.0930]) even while increasing minority exposure overall. This is an important failure case: the reranker is not collapsing uniformly, but it is still uneven across target cultures. The cleaner routeA_small Gemini track shows much more consistent per-target gains, which suggests that the remaining bottleneck is not missing model capacity alone.

Instead, the stronger explanation is data-level confounding. The V4 audit reports weighted source predictability from culture of 0.911765 on V4 main and 1.0 on routeA_small, so source confound remains substantial even after the V4 contract cleanup. We therefore interpret routeA_small as a sanity-check track rather than a primary benchmark, and we treat the remaining per-culture asymmetry on V4 main as a limitation that motivates future PAL-guided repair and stronger source balancing.
